\documentclass[10pt]{beamer}
\usepackage{pgfpages}
\setbeameroption{show notes}
\setbeameroption{show notes on second screen=right}
\usetheme[progressbar=frametitle]{metropolis}
% workaround to set theme back with pdfpages notes
\makeatletter 
\def\beamer@framenotesbegin{% at beginning of slide
     \usebeamercolor[fg]{normal text}
      \gdef\beamer@noteitems{}% 
      \gdef\beamer@notes{}% 
}
\makeatother

\usepackage{natbib}
\usepackage{appendixnumberbeamer}
\usepackage{pbox}
\usepackage{booktabs}
\usepackage[scale=2]{ccicons}
\usepackage{caption}
\usepackage{pgfplots}
\usepgfplotslibrary{dateplot}
\usepackage[]{amsmath}
\usepackage{xspace}
\newcommand{\themename}{\textbf{\textsc{metropolis}}\xspace}
\bibliographystyle{plain}
\title{Nuclear Waste: Too Hot to Touch}
\subtitle{Physics Comps Presentation}
\date{\today}
\date{}
\author{Yuping Huang}
\institute{Department of Physics and Astronomy, Carleton College}
\titlegraphic{\hfill\includegraphics[height=1.5cm]{carleton.pdf}}
\AtBeginSection[]
{
   \begin{frame}{Outline}
  \setbeamertemplate{section in toc}[sections numbered]
       \tableofcontents[currentsection]
   \end{frame}
}
\begin{document}
\begin{frame}
\begin{figure}
\includegraphics[height=\textheight]{last}
\caption*{Image credit XKCD.com}
\end{figure}
\end{frame}
\maketitle
\note[itemize]{
\item thanks and welcome
\item set the scene--some number
}
\begin{frame}{Introduction: why does it matter?}
\begin{itemize}
	\item World nuclear energy consumption will increase by 40\% by 2030
    \item Over \alert{10,000} metric tons of highly radioactive waste generated annually
    \item Highly radioactive waste may remain hazardous for up to millions of years
    \item We will have over \alert{457,000} tons of highly radioactive waste by 2020 without any foreseeable disposal
    scheme
\end{itemize}
\end{frame}
\note[itemize]{
\item 40\% increase energy
\item $10,000$ tons annually
\item harzadous
\item \alert{457,000} tons by 2020
}
\begin{frame}{Spent fuel pool ``near miss'' at Fukushima}
\only<1>{
		\begin{figure}
        \centering
		\includegraphics[width=0.8\textwidth]{vapor}
		\caption*{Water vapor from hydrogen burning in a reactor building at Fukushima}
		\end{figure}
        }
\only<2>{
		\begin{figure}
        \centering
		\includegraphics[width=0.8\textwidth]{exploded}
		\caption*{The explosion in the reactor building}
		\end{figure}
        }
\only<3>{
		\begin{figure}
        \centering
		\includegraphics[width=0.8\textwidth]{fukushima_nocool}
		\caption*{The Japan Times article on the cooling halt of the spent fuel pool}
		\end{figure}
        }
\end{frame}

\begin{frame}{The nightmare scenario}
\note[item]{1 megabecquerel per square meter -- trigger relocation order}
\note[item]{jan, apr, jul, oct, 2015}
\note[item]{convince you of the gravity of this topic}
\begin{figure}
\centering
\only<1>{\includegraphics[width=0.6\textwidth]{jan}\caption*{Jan, 2015}}
\only<2>{\includegraphics[width=0.6\textwidth]{oct}\caption*{Oct, 2015}}
\end{figure}
An estimate of the spread of Cesium-137 radioactive cloud if the spent fuel pools catch fire at the Peach Bottom nuclear power plant in Pennsylvania at different time of the year. \cite{s16}
\end{frame}

\begin{frame}{Outline}
  \setbeamertemplate{section in toc}[sections numbered]
  \tableofcontents[hideallsubsections]
\end{frame}

\section{Radiation and Doses}
\begin{frame}{Unstable nuclei}
  \begin{columns}[T,onlytextwidth]
    \column{0.4\textwidth}
      \begin{itemize}
      	\item Configuration decides stability
      	\item Binding energy = Coulomb + Nuclear
      	\item Minimum energy at $Z\approx N$
      \end{itemize}

    \column{0.6\textwidth}
		\begin{figure}
		\centering
		\includegraphics[height=0.8\textheight]{stability.png}
		\caption*{half-lives of isotopes with different neutron and proton numbers. (Image credit Wikipedia)}
        \end{figure}
  \end{columns}
\end{frame}
\note[itemize]{
\item binding energy as a function of number of protons and neutrons, nuclear force + Coulomb force
\item stable configurations minimize the function
\item as there are more protons. need neutrons to weaken the long-range Coulomb repulsion
}

\begin{frame}{Radioactive decay}
\note[item]{set N=1/2$N_0$ to get half-life}
\note[item]{subsequent decay of daughter nuclei could happen which makes things harder; if half life of daughter much shorter then we can treat it as one exponential; otherwise there will be some equilibrium of daughter nuclei and their subsequent decay products}
\begin{itemize}[<+->]
\item Exponential decay: constant probability of decay $\lambda$ from unstable high
energy nucleus state to lower energy state
\[\lambda = -\frac{(dN(t)/dt)}{N(t)} \Rightarrow N(t)=N_0 e^{-\lambda t}\]
\item Half-life $t_{1/2}=\frac{\ln 2}{\lambda}$: the time required for half of the sample to decay on average
\item{(Radio)activity: number of decays per unit time}
\[\mathcal{A}(t)\equiv \lambda N(t)\]
\item Units: becquerel (Bq) = 1 decay/s or curie (Ci) = $3.7\times 10^{10}$ Bq
\item Specific activity: number of decays per unit time per mass
\end{itemize}
\end{frame}

\begin{frame}{Modes of decay}
\begin{itemize}[<+->]
\item $\alpha$ decay: emits helium-4 atom ${}^4_2$He
\[{}^A_Z X_N \rightarrow {}^{A-4}_{Z-2}X'_{N-2} + \alpha.\]
\item $\beta^-$ decay: emits electrons ($\beta$ radiation) and neutrinos  
\[{}^A_Z X_N \rightarrow {}^A_{Z+1} X'_{N+1} + e^- + \bar{\nu}_e\]
\item $\gamma$ decay: normally occurs after other forms of decays for nuclei in excited state,
emits $\gamma$ ray in the range of $0.1$ to $10$ MeV
\end{itemize}
\end{frame}

\begin{frame}{Half-life and specific activity of radioactive isotopes}{fragile}
  \begin{table}
  \centering
    \begin{tabular}{llll}
      \toprule
      Isotope & Half-life & Specific Activity (Ci/g) \\
      \midrule
      Cesium-137 & 37 years & $88$\\
      Radium-226 & 1600 years  & $1$\\
      Plutonium-239 & 24,000 years & $0.063$ \\
      Uranium-233 & 1.6 million years & $9.3\times 10^{-3}$\\
      Thorium-232 & 14 billion years & $1.1\times 10^{-7}$ \\
      \bottomrule
    \end{tabular}
  \end{table}
  NB: The isotope used for computerized tomography (CT) scan has activity $\sim 10^{-2}$ Ci.
\end{frame}
\note[itemize]{
\item cesium-137 most concerning
\item plutonium makes atomic bombs
\item u-232 long half life and not safe level of radiation
\item thorium-232 can go thru neutron capture to form uranium-233
}

\begin{frame}{Health Effects and Dose}
\begin{itemize}[<+->]
\item Radiation deposit energy and ionize biological molecules, e.g. the DNA
\item \emph{Equivalent dose} $H_R$: energy deposited by the radiation per mass weighted by the type of radiation
\[H_R = w_R \times \frac{E_R}{m_b},\]
$E_R$: energy of radiation, $m_b$ mass of recipient, $w_R$: weighting factor of radiation type.
\item Unit: sievert (Sv) = J/kg, equivalent to the effect of $1$J of x-rays on $1$kg of biological mass
\item For $\gamma$ radiations, the equivalent dose rate is approximately
\[\frac{dH_R}{dt}(\mu\text{Sv h}^{-1}) \approx \frac{\mathcal{A}\text{(MBq)}\times E_{\gamma}\text{(MeV)}}{6\times [r\text{(m)}]^2}\]
\end{itemize}
\end{frame}

\begin{frame}{Weighting factor $w_R$}
The equivalent dose $H_R = w_R \times \frac{E_R}{m_b}.$
\begin{table}
\begin{tabular}{l c c}
\toprule
Type of radiation & Energy range & Weighting factor, $w_R$\\
\hline
Photons, electrons & All energies & 1\\
Neutrons & $<10$ keV & 5 \\
& $10-100$ keV & 10 \\
& $100$ keV-$2$ MeV & 20 \\
& $2-20$ MeV & 10 \\
Protons & $<20$ MeV & 5 \\
$\alpha$ particles& & $20$\\
\bottomrule
\end{tabular}
\caption*{Weighting factors $w_R$ for different types of radiations. \cite{icrp74}}
\end{table}
\end{frame}
\begin{frame}{Cesium-137}
\note[item]{Image shown is the city of Priyat in Ukraine which is inhabitable but is now open to tourist visit}
\note[item]{Dose level still 80 times higher than the recommended maximum for external irradiation of the human body}
\begin{figure}
\centering
\includegraphics[width=0.8\textwidth]{Priyat}
\caption*{After the Chernobyl accident in which 27 kg of cesium-137 were released, a 4,000 km${}^2$ area still currently contains too much cesium-137 to be inhabited (with dose rate $1\mu$Sv/h)}
\end{figure}

\end{frame}

\begin{frame}{Long half-life vs short half-life}
\begin{figure}
\centering
\only<1->{
\includegraphics[width=0.6\textwidth]{csur.png}}
\note[item]{0.01 Ci is enough for medical radiation diagnosis like (CT)}
\note[item]{it's not gonna kill you immediately, but it is not safe}
\only<1>{\caption*{Activity of 1 gram of cesium-137 vs uranium-233 over time, where blue indicates cesium-137 and green uranium-233}}
\only<2>{\caption*{\alert{Conclusion: Long half-life isotopes can be tricky to handle}}}
\end{figure}
\end{frame}


\section{Sources of Nuclear Waste}
\begin{frame}{Categories of nuclear waste (defined by the NRC) \cite{s01}}
\begin{itemize}[<+->]
\item \alert<13>{High level waste}
	\begin{itemize}
    \item typically comes from spent fuel of nuclear reactor
	\item generates large amount of heat and radiation
    \item contains isotopes of varying half life
    \end{itemize}
\item Transuranic waste
	\begin{itemize}
    \item materials contaminated with artificially made radioactive isotopes that have atomic number larger
    than uranium
    \item mostly comes from the military
    \item low activity but long half life
    \end{itemize}
\item Low-level waste
	\begin{itemize}
    \item Low activity waste. Examples are papers and tools used in cleanup operations in the reactor.
    \end{itemize}
\item Uranium tailings
	\begin{itemize}
	\item radioactive thorium and radium left in the uranium mine after the uranium is extracted
    \end{itemize}
\end{itemize}
\end{frame}

\note[itemize]{
\item tailing: radium can decay into radon gas which is an alpha-emitter; be careful not to pollute water and air
}

\begin{frame}{What happens to nuclear fuel: neutron capture}
\note{Thermal neutrons means most probable temperature at $17\deg$C}
\begin{itemize}[<+->]
\item Fresh nuclear fuel: 97\% uranium-238 and 3\% uranium-235
\item The fuel rod is irradiated with \alert{thermal ($0.025$ eV)} neutrons to induce neutron absorption
\item Two main outcome: fission and radiative capture
\end{itemize}
\begin{figure}
\centering
\includegraphics[width=0.8\textwidth]{rod}
\end{figure}
\end{frame}

\begin{frame}{Nuclear fission}
\begin{figure}
\centering
\includegraphics[width=0.9\textwidth]{fission}
\caption*{Schematic illustration of the thermal neutron induced fission of uranium-235. \cite{l01}}
\end{figure}

\end{frame}
\note[itemize]{
\item Thermal neutron gets absorbed; the nucleus becomes unstable and split into two parts
\item the two parts repel; KE of neutron converted to vibrational energy
\item as the amplitude gets higher nuclear force potential drops off sharply and coulomb potential
dominates
\item boom. gigantic binding energy difference produces energy.
}

\begin{frame}{Radiative capture}
\begin{itemize}[<+->]
\item For a nucleus that is not fissile, it absorbs a neutron, form a heavier nuclei and radiates to return to ground state
\item uranium-238 in nuclear fuel can generate plutonium-239
\[{}^{238}\text{U} \xrightarrow{\text{radiative capture}} {}^{239}\text{U} \xrightarrow{\text{half-life 23 min}} {}^{239}\text{Np}\xrightarrow{\text{half-life } 2.4 \text{ d}} {}^{239}\text{Pu}\]
\item Thorium-232 also goes through radiative capture to form uranium-233
\end{itemize}
\begin{figure}
\centering
\includegraphics[height=0.5\textheight]{pu}
\end{figure}
\end{frame}

\begin{frame}{Composition of spent fuel}
\begin{itemize}
\item Nuclear fission produces fission products like cesium-137, strontium-90 with half-lives shorter than
40 years
\item Neutron capture produces uranium-233, plutonium and other transuranic elements (minor actinides) with half-lives up to
millions of years
\end{itemize}
\begin{figure}
\centering
\only<1>{\includegraphics[height=0.5\textheight]{pu}}
\only<2>{\includegraphics[height=0.5\textheight]{spent-fuel}}
\only<2>{\caption*{Composition of spent fuel from light-water reactor; exact composition varies by reactor designs}}
\end{figure}
\end{frame}

\section{Nuclear Waste Management}
\begin{frame}{Shielding radiations}
	\note[item]{SLOW DOWN}
    \note[item]{Interaction probability is independent of distance traveled}
    \note[item]{Decay probability is independent of the time passed}
	\begin{itemize}[<+->]
    	\item For spent fuel, radiations and neutrons need to be shielded
    	\item Shielding forces the radiations to interact with the shielding media
        \item Nuclear cross section $\sigma$ quantifies the effective interaction probability
        \item The interaction probability is independent of the distance traveled
        \item Intensity $I(z)$ of incident beam falls off exponentially as thickness $z$ increases
        \[I(z)=I(0) e^{-\mu z}, \mu \propto \sigma.\]
    \end{itemize}
\end{frame}

\begin{frame}{Nuclear cross section\ footnote{section adapted from \cite{k88}}}
	\note[item]{erase the board and do dimensional analysis}
	\note[item]{tells you how many reactions can happen with one nucleus}
    \note[item]{make the point about 10 times as thick water}
	\begin{itemize}[<+->]
	\item Nuclear cross section characterizes reaction probability and has the unit of area per nucleus
    \[\sigma = \frac{R}{I_i N},\]
    where $I_i$ is incident beam intensity, $R$ is the number of reaction occurred per time, $N$ the number of nuclei in the target
	\item Typical nuclei radii are $\sim 10^{-14}$ and have effective area $\pi r^2\sim 10^{-28}=1$ barn
    \item Boron nucleus neutron capture has $\sigma = 3840$ barns
    \item $\gamma$ ray interaction with water has $\sigma \sim 0.01$ barns, with concrete $\sigma\sim 0.1$ barns
    \end{itemize}
\end{frame}

\begin{frame}{Spent fuel pools}
\note[item]{heat drops by more than 90\% for the first year}
\note[item]{Two years is enough but there aren't enough subsequent storage space available}
  \begin{columns}[T,onlytextwidth]
    \column{0.5\textwidth}
      \begin{itemize}[<+->]
      	\item Removed spent fuel typically spend 5 years or more in spent fuel pools
        \item Filled with boron-rich water because boron atoms have a larger neutron capture cross-section
        \item The water is circulated and cooled at $50^{\circ}$C
        \item At least 15-meter deep to shield gamma rays
      \end{itemize}

    \column{0.5\textwidth}
		\begin{figure}
		\includegraphics[width=0.9\textwidth]{fukushimafuel-rods}
		\caption*{Spent fuel pool at Fukushima}
		\end{figure}
  \end{columns}
\end{frame}

\begin{frame}{Dry cask storage}

  \begin{columns}[T,onlytextwidth]
    \column{0.5\textwidth}
      \begin{itemize}[<+->]
      	\item After 5 years in the pool, the spent fuel can be transfered to dry cask storage
        \item Dry casks are sturdy metal and concrete structures designed to withstand harsh conditions
        \item Air is circulated to cool the cask
        \item The NRC believes that dry casks storage will be safe for at least 100 years
      \end{itemize}

    \column{0.5\textwidth}
		\begin{figure}
		\includegraphics[width=0.9\textwidth]{dry_storage.jpg}
		\caption*{A dry cask storage stie}
		\end{figure}
  \end{columns}
\end{frame}
\section{Nuclear Waste Disposal}
	\begin{frame}{Deep Geological Repository}
    \begin{itemize}[<+->]
    \item The goal of nuclear waste disposal is to minimize radiation exposure for over $10^6$ years and to prevent proliferation
    \item The science community generally accepts that deep (over $300$m) underground disposal for high-level waste is much safer than above-ground storage
    \item Other alternatives include transmutation and space storage, which are too costly
    \end{itemize}
    \end{frame}
    
    \begin{frame}{Yucca Mountain Nuclear Waste Repository}
    \note[item]{established Nuclear Waste Policy Act of 1987}
    \note[item]{later became the only site studied for feasibility}
    \note[item]{the Department of Energy was directed not to develop complementary approaches}
    \note[item]{risks: oxidation from moise air; water percolation to water table; hot upwelling water from below}
    \only<1>{
    \begin{figure}
    \centering
    \includegraphics[width=\textwidth]{yuccaRepo}
    \caption*{Scheme of the Yucca Mountain Nuclear Waste Repository; the depth of the repository is about $300$ meters}
    \end{figure}
    }
    \only<2->{
    \note[item]{prevent oxidation is the responsibility of the packaging of the waste, but it is hard to predict its integrity for over 1 million years}
    Risk factors:\\
    \begin{itemize}
    	\item Oxidation of the waste from moist air
        \item Water percolation from the surface to the water table
        \item \alert<2>{Tectonic motion and climate change}
        \item Hot upwelling water from below
        \item \alert<3>{Catastrophic events}
    \end{itemize}
    }
    \end{frame}
	\begin{frame}{Volcanism: risk assessment}
    \note[item]{disruptive volcanic eruption releasing radioactive material into the air}
    \note[item]{observation and inferences; human choices; computational models}
    \centering
        \begin{figure}
		\includegraphics[height=0.6\textheight]{volcano}
        \caption*{A flowchart for the volcanism risk estimation problem. The oval-shaped parts are observations and statistical inferences, the rectangular ones human choices, and the hexagon the DOE computational model. \cite{me06}}
		\end{figure}
    \end{frame}
	\begin{frame}{How to deal with uncertainty?}
    	\only<1>{
        \begin{figure}
		\includegraphics[height=0.8\textheight]{meme}
		\end{figure}
        }
        
        \only<2>{
    	\begin{figure}
		\includegraphics[height=0.8\textheight]{bill}
		\caption*{Propagation of uncertainty. Courtesy of Bill Titus}
		\end{figure}
        }
        \only<3>{
        \alert{BUT}
        \begin{itemize}
        \item We don't necessarily have an analytical expression for the model
        \item The linear approximation might fail if the uncertainties of the input parameters
        are large
        \item The error bar may not contain all the information about the uncertainty
        \end{itemize}
        }
    \end{frame}
    
    \begin{frame}{Monte Carlo simulations}
    	Suppose we have a function $y(x_1,x_2,...,x_n)$ that is not necessarily analytic (for example, $y$ can be the expected leaked dose calculated from a simulation with a bunch of input parameters).\\
        \begin{itemize}
            \item Generate a large number of $(x_1,x_2,...x_n)$ samples according to their respective distributions.
            \item For each $(x_1,x_2,...,x_n)$ sample, calculate a $y$ value (with the computational model).
            \item Aggregate all the $y$'s generated to estimate its distribution or uncertainty.
        \end{itemize}
    \end{frame}
    
%    \begin{frame}{Monte Carlo simulations: an example}
%		\begin{figure}
%		\includegraphics[width=0.9\textwidth]{Monte-Carlo-laundry-example}
%		\caption*{An example of Monte Carlo simulation to infer environmental impact of laundry based
%        on a user's laundry habit}
%		\end{figure}
%    \end{frame}
    \begin{frame}{Monte Carlo simulation: an example}{fragile}
    \only<1>{
        \alert{You are in a lab.}
    	\begin{itemize}
        \item Task: find out how many nuclei will remain after $t=200$s for the sample you have
        \item You \alert{don't} know the exponential decay law $N(t)=N_0 e^{-\lambda t}$
        \item You have: a program \textsc{NOTNDSOLVE} takes in $N_0$, $\lambda$, $t$ and predicts $N(t)$
        \item Measurements made: $N_0= (1\pm 0.2)\times 10^{27}$, $\lambda = 0.003\pm 0.001$/s
        \item Assume each parameter is uniformly distributed inside its error bar 
        \end{itemize}
      }
      
      \only<2>{
      \alert{Here's what you do}\\
      Repeat the following $10,000$ times
      \begin{itemize}
      \item Generate a random $N_0$ and a random $\lambda$ from their respective distributions
      \item Calculate $N(t=20s)$ using $NOTNDSOLVE$
      \end{itemize}
      Now plot a histogram of all the $N(t=20s)$ numbers you have. It quantifies the uncertainty of your prediction
      }
      
      \only<3->{
      Probability distribution for the predicted $N(t=200s)$
      \begin{figure}
      \centering
      \only<3>{
      \includegraphics[height=0.7\textheight]{nList}
      \caption*{$10^4$ samples;Mean: $5.52\times 10^{26}$; standard deviation $5.53\times 10^{26}$}
      }
      \only<4>{
      \includegraphics[height=0.7\textheight]{more}
      \caption*{$10^6$ samples; Mean: $5.52\times 10^{26}$; standard deviation $9.04\times 10^{25}$}
      }
      \end{figure}
      
      }
    \end{frame}
    
    \begin{frame}{Volcanic disruption problem revisited}
        \begin{figure}
		\includegraphics[height=0.6\textheight]{volcano}
        \caption*{A flowchart for the volcanism risk estimation problem. The oval-shaped parts are observations and statistical inferences, the rectangular ones human choices, and the hexagon the DOE computational model. \cite{me06}}
		\end{figure}
    \end{frame}
    
    \begin{frame}{The volcanism probability: NRC vs DOE}
    \note[item]{on the edge situation}
    \note[item]{if the expected leaked dose is really high, then the risk is higher than when the leaked dose is low}
    \begin{itemize}[<+->]
    	\item EPA guideline: if the disruptive event rate is lower than \alert{$10^{-8}$} events per year then the risk is negligible
    	\item The DOE estimates an disruptive volcanic eruption rate of $1.5\times 10^{-9}$ events per year
        with 90 percent confidence bound \alert{$5.4\times 10^{-9}$-$4.9\times 10^{-8}$} events per year
        
        \item The NRC uses a subset of the DOE's models that reflect conservatism and concludes a disruption rate of \alert{$10^{-7}$ - $10^{-8}$} events per year
        
        \item In an edge case like this, we need to understand better the effects of the eruption on the leaked dose to assess the risk
    \end{itemize}
    \end{frame}

\begin{frame}{Status of Yucca Mountain}
\note[item]{Obama administration defunded the project in 2011}
Steven Chu, former secretary of energy said in 2009 \cite{aa12}:\\
\vspace{0.5cm}
``Yucca Mountain as a repository is off the table. What we’re going to be doing is saying, let’s step back. We realize that we know a lot more today than we did 25 or 30 years ago...(dry cask storage) gives us time to figure out what we should do for a long-term strategy.''\\
\vspace{0.5cm}
In 2011, the project was defunded by the Obama administration
\end{frame}

\begin{frame}{Newsflash}
\only<1>{
\begin{figure}
\centering
\caption*{Las Vegas Review Journal, Jan 19, 2017}
\includegraphics[width=0.8\textwidth]{rickperry.png}
\end{figure}
}
\only<2>{
\begin{figure}
\centering
\caption*{Las Vegas Review Journal, Jan 31, 2017}
\includegraphics[width=0.8\textwidth]{fuck.png}
\end{figure}
}
\end{frame}

\section{Discussions}

\begin{frame}{Discussions: what else can be done?}
\begin{itemize}[<+->]
\item Deep underground storage is complicated, both politically and scientifically
\item Dry cask storage can buy the US a few decades to come up with different solutions (e.g. fast neutron reactors)
\item France is building a reversible long-term ($\sim 100$ years) underground storage
\item Sweden and Finland started constructing deep geological repositories
\end{itemize}

\end{frame}

\begin{frame}[allowframebreaks]{References}
  \bibliography{refs}

\end{frame}

\begin{frame}{Acknowledgment}
\begin{itemize}
\item Thanks to you for coming
\item Thanks to Yuhao Wan, Sam McHugh, and Gaston for attending the practice talk
\item Thanks to advisors Joel Weisberg and Alex McMurty
\end{itemize}
\end{frame}

\begin{frame}[standout]
\usebeamercolor[bg]{normal text}
Questions?
\end{frame}

\end{document}
